% Practice sheet — 2 columns, topic-wise, answer under every problem, full solutions at the end
% Switch \solutionsatend <-> \solutionsinline (ask the student which they want)
\documentclass[10pt]{article}
\usepackage[tiro]{edtech}
\geometry{left=12mm,right=12mm,top=12mm,bottom=20mm,footskip=10mm}
\etset{subject={পদার্থবিজ্ঞান ১ম পত্র}, chapter={অধ্যায় ০৩: গতিবিদ্যা}}
\pagestyle{barfoot}
\setlength{\columnseprule}{0.4pt}\setlength{\columnsep}{6mm}
\solutionsatend
\begin{document}
\begin{center}{\headfont\bfseries\Large অধ্যায়-৩}\\[-2pt]{\headfont\bfseries\LARGE গতিবিদ্যা}\end{center}
\begin{multicols}{2}
\topicbar{Topic-1: সমত্বরণ ও মুক্তভাবে পড়ন্ত বস্তু}
\problem{একটি গাড়ি স্থির অবস্থান থেকে \SI{2}{m.s^{-2}} সমত্বরণে চলতে শুরু করে। \SI{10}{s} পরে গাড়িটির বেগ ও অতিক্রান্ত দূরত্ব নির্ণয় কর।}[বোর্ড স্টাইল]{\SI{20}{m.s^{-1}}, \SI{100}{m}}
\solution{$v=u+at=0+2\times10=\SI{20}{m.s^{-1}}$\\ $s=ut+\frac12at^2=0+\frac12\times2\times10^2=\SI{100}{m}$}

\problem{একটি বস্তুকে \SI{49}{m.s^{-1}} বেগে খাড়া উপরের দিকে নিক্ষেপ করা হলো। বস্তুটির সর্বোচ্চ উচ্চতা ও উত্থানকাল নির্ণয় কর। [$g=\SI{9.8}{m.s^{-2}}$]}[বোর্ড স্টাইল]{\SI{122.5}{m}, \SI{5}{s}}
\solution{সর্বোচ্চ উচ্চতায় $v=0$।\\ $H=\dfrac{u^2}{2g}=\dfrac{49^2}{2\times9.8}=\SI{122.5}{m}$\\ $t=\dfrac{u}{g}=\dfrac{49}{9.8}=\SI{5}{s}$}

\problem{সুষম ত্বরণে চলমান একটি বস্তু পঞ্চম সেকেন্ডে \SI{18}{m} এবং অষ্টম সেকেন্ডে \SI{27}{m} দূরত্ব অতিক্রম করে। বস্তুটির আদিবেগ ও ত্বরণ নির্ণয় কর।}[ভর্তি পরীক্ষা স্টাইল]{$u=\SI{4.5}{m.s^{-1}}$, $a=\SI{3}{m.s^{-2}}$}
\solution{$n$-তম সেকেন্ডে অতিক্রান্ত দূরত্ব $s_n=u+\frac12a(2n-1)$\\
$18=u+4.5a$ \dots (i)\qquad $27=u+7.5a$ \dots (ii)\\
(ii) $-$ (i): $3a=9\Rightarrow a=\SI{3}{m.s^{-2}}$\\
(i) হতে: $u=18-4.5\times3=\SI{4.5}{m.s^{-1}}$}

\problem{একটি ট্রেন \SI{72}{km.h^{-1}} বেগে চলছিল। ব্রেক চাপার পর ট্রেনটি \SI{200}{m} দূরত্ব অতিক্রম করে থেমে গেল। ট্রেনটির মন্দন নির্ণয় কর।}[বোর্ড স্টাইল]{\SI{1}{m.s^{-2}}}
\solution{$u=\SI{72}{km.h^{-1}}=\dfrac{72\times1000}{3600}=\SI{20}{m.s^{-1}}$,\ $v=0$\\ $v^2=u^2-2as\Rightarrow a=\dfrac{u^2}{2s}=\dfrac{400}{400}=\SI{1}{m.s^{-2}}$}

\topicbar{Topic-2: প্রক্ষেপক ও বৃত্তাকার গতি}
\problem{একটি প্রক্ষেপককে অনুভূমিকের সাথে \SI{30}{\degree} কোণে \SI{40}{m.s^{-1}} বেগে নিক্ষেপ করা হলো। এর অনুভূমিক পাল্লা নির্ণয় কর। [$g=\SI{9.8}{m.s^{-2}}$]}[বোর্ড স্টাইল]{\SI{141.4}{m}}
\solution{$R=\dfrac{u^2\sin2\theta}{g}=\dfrac{40^2\times\sin60^\circ}{9.8}=\dfrac{1600\times0.866}{9.8}=\SI{141.4}{m}$}

\problem{একটি বস্তুকে অনুভূমিকের সাথে \SI{30}{\degree} কোণে \SI{98}{m.s^{-1}} বেগে নিক্ষেপ করা হলো। বস্তুটির বিচরণকাল নির্ণয় কর।}[অনুশীলনী]{\SI{10}{s}}
\solution{$T=\dfrac{2u\sin\theta}{g}=\dfrac{2\times98\times\sin30^\circ}{9.8}=\dfrac{98}{9.8}=\SI{10}{s}$}

\problem{একটি টাওয়ারের শীর্ষ থেকে একটি বল \SI{20}{m.s^{-1}} বেগে অনুভূমিকভাবে নিক্ষেপ করা হলো। বলটি \SI{3}{s} পর ভূমিতে পড়লে টাওয়ারের উচ্চতা কত এবং বলটি টাওয়ারের পাদদেশ থেকে কত দূরে পড়বে?}[অনুশীলনী]{\SI{44.1}{m}, \SI{60}{m}}
\solution{উল্লম্ব দিকে আদিবেগ শূন্য।\\ $h=\frac12gt^2=\frac12\times9.8\times3^2=\SI{44.1}{m}$\\ $x=u\,t=20\times3=\SI{60}{m}$\\[3pt]
\begin{tikzpicture}[scale=0.42,font=\scriptsize]
 \fill[black!15](0,0) rectangle (0.8,4.41); \draw(0,0) rectangle (0.8,4.41);
 \draw[thick,domain=0:6,smooth,variable=\x,etblue] plot ({0.8+\x},{4.41-4.41*(\x/6)^2});
 \draw(-0.5,0)--(7.5,0); \draw[{Latex}-{Latex}](0.8,-0.5)--(6.8,-0.5) node[midway,below]{\SI{60}{m}};
 \draw[{Latex}-{Latex}](-0.3,0)--(-0.3,4.41) node[midway,left]{\SI{44.1}{m}};
 \draw[-{Latex}](0.8,4.41)--(2.4,4.41) node[above]{\SI{20}{m.s^{-1}}};
\end{tikzpicture}}

\problem{\SI{0.5}{m} ব্যাসার্ধের বৃত্তাকার পথে একটি কণা প্রতি মিনিটে \num{120} বার ঘোরে। কণাটির কেন্দ্রমুখী ত্বরণ নির্ণয় কর।}[অনুশীলনী]{\SI{78.96}{m.s^{-2}}}
\solution{$f=\dfrac{120}{60}=\SI{2}{Hz}$,\quad $\omega=2\pi f=4\pi=\SI{12.57}{rad.s^{-1}}$\\ $a=\omega^2r=(12.57)^2\times0.5=\SI{78.96}{m.s^{-2}}$}

\topicbar{Topic-3: MCQ অনুশীলন}
\pmcq[ধারণাগত]{বেগ–সময় লেখচিত্রের ঢাল নির্দেশ করে—}{সরণ}{ত্বরণ}{দ্রুতি}{বল}{(খ) ত্বরণ}
\pmcq[ধারণাগত]{প্রক্ষেপকের অনুভূমিক পাল্লা সর্বাধিক হয় যখন নিক্ষেপ কোণ—}{\SI{30}{\degree}}{\SI{45}{\degree}}{\SI{60}{\degree}}{\SI{90}{\degree}}{(খ) \SI{45}{\degree}}
\pmcq[বোর্ড স্টাইল]{স্থির অবস্থান থেকে মুক্তভাবে পড়ন্ত বস্তু প্রথম, দ্বিতীয় ও তৃতীয় সেকেন্ডে যে দূরত্ব অতিক্রম করে তাদের অনুপাত—}{$1:2:3$}{$1:4:9$}{$1:3:5$}{$1:1:1$}{(গ) $1:3:5$}
\pmcq[ধারণাগত]{সর্বোচ্চ উচ্চতায় প্রক্ষেপকের বেগ কত?}{শূন্য}{$u\sin\theta$}{$u$}{$u\cos\theta$}{(ঘ) $u\cos\theta$}
\pmcq[বোর্ড স্টাইল]{কোনো কণার সরণ $x=3t^2+2t$ (m) হলে $t=\SI{2}{s}$ সময়ে কণাটির বেগ কত?}{\SI{12}{m.s^{-1}}}{\SI{14}{m.s^{-1}}}{\SI{16}{m.s^{-1}}}{\SI{10}{m.s^{-1}}}{(খ) \SI{14}{m.s^{-1}}\quad[$v=\frac{dx}{dt}=6t+2$]}
\pmcq[ধারণাগত]{কৌণিক বেগের SI একক কোনটি?}{\si{m.s^{-1}}}{\si{rad.s^{-2}}}{\si{rad.s^{-1}}}{\si{Hz.m}}{(গ) \si{rad.s^{-1}}}
\end{multicols}

\printsolutions[সমাধান — Topic-1 \& Topic-2]
\end{document}
